\section{Canonical count phases and complete-tail errors}\label{sec:highcount}

The two results in this section use the actual canonical source tail to
quantify the analytical count record. Their proof notes
\cite{TNDHighCount,TNDCountTail} were accepted within scope by separate
same-model internal LLM assessments
\cite{TNDHighCountReview,TNDCountTailReview}. This supplies no specialist
independence or new physical access after nonzero counts. We give the
complete scalar proofs rather than infer a growing-order conclusion from
fixed-order domain estimates.

\subsection{Complete high-count concentration}
For $r\ge0$, $x=\pi\ee^{4r}$, every term of the full source sum is positive.
In addition to Lemma~\ref{lem:source}, the explicit comparison is
\begin{equation}\label{eq:hccompare}
 \begin{gathered}
 c_-\ee^{9r-\pi\ee^{4r}}\le\Phi(r)\le c_+\ee^{9r-\pi\ee^{4r}},\\
 c_-=\pi^2(2-3/\pi),\qquad
 c_+=2\pi^2\sum_{n\ge1}n^4\ee^{-\pi(n^2-1)}<\infty.
 \end{gathered}
\end{equation}
The first term gives only the positive lower bound. The upper bound pays
all $n\ge1$ using $\ee^{-(n^2-1)x}\le\ee^{-\pi(n^2-1)}$; no finite
prefix replaces the source.

Define, for $j\ge1$,
\begin{equation}\label{eq:hccenter}
 \begin{gathered}
 j=2\pi a_j\ee^{4a_j},\qquad M_j=j/a_j,\qquad s_j=M_j^{-1/2},\\
 \nu_{j,\lambda}(\dd r)=
 \frac{r^{2j}\ee^{\lambda r^2}\Phi(r)\dd r}
 {\int_0^\infty r^{2j}\ee^{\lambda r^2}\Phi(r)\dd r}.
 \end{gathered}
\end{equation}
The center is unique and positive. The complete count integral
\eqref{eq:dj}, including its original signed amplitudes and exact angular
trace one, gives for $t>0$
\begin{equation}\label{eq:hcratio}
 \frac{d_j(t,k)}{P_j(t)}=\E_{\nu_{j,\lambda(t)}}\cos(kr).
\end{equation}
The full-line factors of two and the positive count prefactor cancel;
the absolute calibration of $d_j$ remains unchanged.

\begin{theorem}[Uniform concentration and count phases]\label{thm:highcount}
There are finite constants $C_1,C_2$ and an integer $J$ such that, for
$j\ge J$,
\begin{equation}\label{eq:hcmoments}
 \begin{aligned}
 \sup_{|\lambda|\le1/2}\E_{\nu_{j,\lambda}}|r-a_j|
 &\le C_1\sqrt{a_j/j},\\
 \sup_{|\lambda|\le1/2}\E_{\nu_{j,\lambda}}(r-a_j)^2
 &\le C_2a_j/j.
 \end{aligned}
\end{equation}
For every finite $K$, real $|k|\le K$ and $0<t\le T$,
\begin{equation}\label{eq:hcphase}
 \left|\frac{d_j(t,k)}{P_j(t)}-\cos(ka_j)\right|
 \le KC_1\sqrt{a_j/j}\longrightarrow0.
\end{equation}
For each fixed real $k\ne0$ and every $0<t\le T$,
\begin{equation}\label{eq:hcsigns}
 \limsup_{j\to\infty}\frac{d_j(t,k)}{P_j(t)}=1,\qquad
 \liminf_{j\to\infty}\frac{d_j(t,k)}{P_j(t)}=-1.
\end{equation}
Both signs therefore occur infinitely often in the literal $d_j(t,k)$.
The two realizing subsequences can be chosen independently of $t$.
\end{theorem}
\begin{proof}
Abbreviate $a=a_j$, $M=j/a$, $s=M^{-1/2}$ and put
$q_j(r)=2j\log r-\pi\ee^{4r}$,
$f_{j,\lambda}(r)=q_j(r)+\lambda r^2+9r$. For $h=r-a>-a$, the center
identity gives exactly
\begin{equation}\label{eq:hcexponent}
 q_j(a+h)-q_j(a)
 =2j[\log(1+h/a)-h/a]-\frac M2[\ee^{4h}-1-4h].
\end{equation}
The logarithmic bracket is nonpositive. With $c_0=4\ee^{-4}$, Taylor's
integral remainder and $\ee^{-4u}-1+4u\ge3u$ for $u\ge1$ imply
\begin{equation}\label{eq:hcthreeparts}
 q_j(a+h)-q_j(a)\le
 \begin{cases}
 -4Mh^2,&h\ge0,\\
 -c_0Mh^2,&-1\le h\le0,\\
 -3M|h|/2,&-a<h\le-1.
 \end{cases}
\end{equation}
The origin is included by its limiting zero density for $j\ge1$.
Since $M=2\pi\ee^{4a}$ and $a\to\infty$, choose $J$ so that
$a\ge1$ and $M\ge(a+9)^2/c_0$ for every $j\ge J$.
This implies $M\ge16,2/c_0,2a,a^2$ and $s\le1/4$.
For $h\ge-1$, uniformly for $|\lambda|\le1/2$, the perturbation is at
most $(a+9)|h|+h^2/2$. The inequalities
\[
 (a+9)|h|\le\frac{c_0M}{4}h^2+\frac{(a+9)^2}{c_0M},\qquad
 \frac{h^2}{2}\le\frac{c_0M}{4}h^2
\]
give, with $c_1=c_0/2$,
\begin{equation}\label{eq:hcglobalupper}
 f_{j,\lambda}(a+h)-f_{j,\lambda}(a)\le1-c_1Mh^2
 \quad(h\ge-1).
\end{equation}
On the entire growing far-left interval $h=-u$, $1\le u<a$, the
perturbation is at most $au$, so $M\ge2a$ yields
\begin{equation}\label{eq:hcleft}
 f_{j,\lambda}(a-u)-f_{j,\lambda}(a)\le-Mu.
\end{equation}
Thus both the complete left tail and the unbounded right tail are paid.

For a positive normalizer, restrict to $0\le h\le s$.
Here $\log(1+y)-y\ge-y^2/2$ for $y\ge0$ and
$\ee^{4h}-1-4h\le8\ee h^2$. The two losses in
\eqref{eq:hcexponent} are at most $1$ and $4\ee$, and the perturbation
loss is at most $ah+h^2/2\le3/2$. With $D_0=3+4\ee$, the normalizing
integral is therefore at least
\begin{equation}\label{eq:hcnormalizer}
 c_-\ee^{f_{j,\lambda}(a)}\ee^{-D_0}M^{-1/2}.
\end{equation}
For $p=1,2$, put $A_p=\int_\R|z|^p\ee^{-c_1z^2}\dd z<\infty$.
Splitting the whole centered-moment numerator at $h=-1$ gives
\begin{align*}
 &\int_0^\infty |r-a|^p r^{2j}\ee^{\lambda r^2}\Phi(r)\dd r\\
 &\quad\le c_+\ee^{f_{j,\lambda}(a)}
 \left[\ee A_pM^{-(p+1)/2}+\int_1^a u^p\ee^{-Mu}\dd u\right]\\
 &\quad\le c_+\ee^{f_{j,\lambda}(a)}
 [\ee A_p+p!]M^{-(p+1)/2}.
\end{align*}
The last left-tail bound extends to $[0,\infty)$ and uses the exact
integral $p!M^{-(p+1)}$. Dividing by \eqref{eq:hcnormalizer} proves
\eqref{eq:hcmoments}, for example with
$C_p=(c_+/c_-)\ee^{D_0}(\ee A_p+p!)$. These constants are independent
of $j,\lambda,t,k,\gamma$ in the stated ranges. This is a separate
growing-$j$ proof using only the zero-order comparison
\eqref{eq:hccompare}.

The real cosine is Lipschitz with constant $|k|$ in $r$, proving
\eqref{eq:hcphase}. The inverse center function is increasing and
unbounded, with
\[
 \frac{da}{dj}=\frac{a}{j(1+4a)}<\frac1{4j},\qquad
 0<a_{j+1}-a_j<\frac14\log(1+1/j)\le\frac1{4j}\to0.
\]
For fixed $k\ne0$, choose first-crossing indices of the levels
$2\pi m/|k|$ and $(2m+1)\pi/|k|$. Their overshoots tend to zero,
so the corresponding cosines tend to $1$ and $-1$; the indices are
eventually distinct. Uniformity in $\lambda$ in \eqref{eq:hcphase}
gives the two subsequence limits simultaneously for all physical $t>0$.
The bound $|d_j/P_j|\le1$ proves the exact limsup and liminf.
Finally $P_j(t)>0$ for $t>0$ transfers the strict signs to $d_j$.
\end{proof}

For each finite $j$, the differentiated source integrals give the exact
signed even-derivative translation
\begin{equation}\label{eq:hcderivatives}
 \begin{aligned}
 \frac{(-1)^j\partial_k^{2j}H_\lambda(k)}
 {\partial_\lambda^jH_\lambda(0)}
 &=\E_{\nu_{j,\lambda}}\cos(kr)\\
 &=\cos(ka_j)+O_K(\sqrt{a_j/j}).
 \end{aligned}
\end{equation}
The denominator is positive. At each fixed $\lambda$ in the slab and
real $k\ne0$ the ratio has the same limsup and liminf. No derivative-order
series, zero theorem or multiplicity claim is inferred.
For every fixed count the radial output is a scalar multiple of
$v^jf_{\lambda(t)}$, hence is Schwartz; \eqref{eq:fullmoment} and
\eqref{eq:radialtail} pay all fixed full-spectrum domains. Their constants
are not uniform in growing $j$.

The center in \eqref{eq:hccenter} was already used in the order-three
Mellin response calculation \cite[Section 2, (3)]{MC03ScalarSaddle}, with
Gaussian scale $\sigma^2=a^2/[2j(1+4a)]$. Its proved limit is on fixed
ordered boxes of a signed three-particle determinant. The present positive
scalar proof on the complete radial domain does not pay its signed
full-chamber allowances or the higher-order obligations discussed in
\cite{MC03HigherOrder}. No new saddle or novelty claim is made.

\subsection{An explicit complete count-tail envelope}
Write $\alpha=\gamma t\in[0,1]$, $\lambda=1/2-\alpha$ and
$h_*=H_{1/2}(0)$. Define positive scalar envelopes
\begin{equation}\label{eq:tailB}
 \begin{gathered}
 I_m(\lambda)=\int_0^\infty r^m\ee^{\lambda r^2}\Phi(r)\dd r,\\
 B_j^{(q)}=\frac{\alpha^j}{j!h_*}I_{2j+q}(\lambda),\qquad q=0,2.
 \end{gathered}
\end{equation}
Thus $B_j^{(0)}=P_j$; $B_j^{(2)}$ is a moment envelope, not an additional
physical readout. The original zero-count convention applies at $\alpha=0$.

First justify complete summation and first differentiation independently.
With the half-line initial probability and full Poisson kernel
\[
 \dd\mu_*(r)=h_*^{-1}\ee^{r^2/2}\Phi(r)\dd r,\qquad
 p_j(y)=\ee^{-y}y^j/j!,
\]
where $p_0(0)=1$, $p_{j>0}(0)=0$, the exact count integrals are
\begin{equation}\label{eq:tailPoisson}
 d_j=\int p_j(\alpha r^2)\cos(kr)\dd\mu_*,\qquad
 B_j^{(q)}=\int r^qp_j(\alpha r^2)\dd\mu_*.
\end{equation}
The complete source bound gives finite initial moments of orders $0,1,2$.
The positive exponential series and Tonelli give
$\sum_jB_j^{(0)}=1$ and $\sum_jB_j^{(2)}=\int r^2\dd\mu_*$.
For $N\ge0$, let
\[
 \tau_N(y)=1-\ee^{-y}\sum_{j=0}^Ny^j/j!\in[0,1].
\]
Its finite-expression derivative is
$\tau_N'(y)=\ee^{-y}y^N/N!\ge0$. Hence, for $q=0,1,2$,
\begin{equation}\label{eq:tailuniformdomination}
 \sup_{0\le\alpha\le1}\int r^q\tau_N(\alpha r^2)\dd\mu_*
 \le\int r^q\tau_N(r^2)\dd\mu_*\longrightarrow0.
\end{equation}
Dominated convergence pays the entire radial domain. Each first real-$k$
derivative has the integrable envelope $rp_j$, and its omitted sum has
envelope $r\tau_N$. Thus the value and first-derivative series converge
uniformly for all real $k$ and the closed alpha slab, and
\begin{equation}\label{eq:tailcomplete}
 \sum_{j\ge0}d_j=\frac{H_{1/2}(k)}{h_*},\qquad
 \partial_k\sum_{j\ge0}d_j=\sum_{j\ge0}\partial_kd_j
 =\frac{H_{1/2}'(k)}{h_*}.
\end{equation}
These passages use their own full-source domination, not a formal
differentiation of a previously established count identity.

\begin{theorem}[Absolute omitted-count errors]\label{thm:counttail}
Let $C_2,J$ be as in Theorem~\ref{thm:highcount} and put
$D=(1+\sqrt{C_2})^2$. Choose a finite integer $N_0\ge1$ satisfying
\begin{equation}\label{eq:tailthreshold}
 N_0+1\ge J,\qquad a_{N_0+1}\ge1,\qquad
 D\frac{a_{N_0+1}^2}{N_0+1}\le\frac12.
\end{equation}
For $N\ge N_0$ define
\begin{equation}\label{eq:tailQ}
 \begin{gathered}
 Q_N(t)=\alpha D\frac{a_{N+1}^2}{N+1}\le\frac{\alpha}{2}<1,\\
 Q_N(t)\le\frac{\alpha D\log^2(N+1)}{16(N+1)}.
 \end{gathered}
\end{equation}
Then, simultaneously for $q=0,2$ and every $j>N$,
\begin{equation}\label{eq:tailratio}
 B_{j+1}^{(q)}\le Q_NB_j^{(q)}.
\end{equation}
For $R_N(t,k)=\sum_{j>N}d_j(t,k)$, every real $k$ and $0\le t\le T$,
\begin{equation}\label{eq:tailerrors}
 |R_N|\le\frac{B_{N+1}^{(0)}}{1-Q_N},\qquad
 |\partial_kR_N|\le
 \frac{\sqrt{B_{N+1}^{(0)}B_{N+1}^{(2)}}}{1-Q_N}.
\end{equation}
\end{theorem}
\begin{proof}
Minkowski and \eqref{eq:hcmoments} yield
\[
 \E_{\nu_{j,\lambda}}r^2
 \le\left(a_j+\sqrt{C_2a_j/j}\right)^2\le Da_j^2
 \quad(j\ge J,\ a_j\ge1).
\]
For $\alpha>0$, the exact adjacent ratios are
\begin{equation}\label{eq:tailadjacent}
 \begin{aligned}
 \frac{B_{j+1}^{(0)}}{B_j^{(0)}}
 &=\frac\alpha{j+1}\E_{\nu_{j,\lambda}}r^2,\\
 \frac{B_{j+1}^{(2)}}{B_j^{(2)}}
 &=\frac\alpha{j+1}\E_{\nu_{j+1,\lambda}}r^2.
 \end{aligned}
\end{equation}
The index $j+1$ in the second ratio is necessary. The exact relation
$a_j^2/j=a_j\ee^{-4a_j}/(2\pi)$ is decreasing for $a_j\ge1$ and tends
to zero. This proves existence of \eqref{eq:tailthreshold} and bounds
both ratios by $Q_N$ for all $j>N$. Since
$4a_m=\log m-\log(2\pi a_m)\le\log m$ for $a_m\ge1$, it also proves
the logarithmic bound in \eqref{eq:tailQ}, namely
$Q_N=O(\alpha\log^2N/N)$. At $\alpha=0$ both positive-index masses
vanish; \eqref{eq:tailratio} holds without using an undefined ratio.

Geometric summation from the correct first seed $N+1$ gives
\[
 T_N^{(q)}:=\sum_{j>N}B_j^{(q)}
 \le B_{N+1}^{(q)}/(1-Q_N),\qquad q=0,2.
\]
Independently justified by \eqref{eq:tailuniformdomination},
\[
 R_N=\int\tau_N(\alpha r^2)\cos(kr)\dd\mu_*,\qquad
 \partial_kR_N=-\int r\tau_N(\alpha r^2)\sin(kr)\dd\mu_*.
\]
The first integral is bounded by $T_N^{(0)}$. Cauchy--Schwarz on the
complete positive tail measure $\tau_N(\alpha r^2)\dd\mu_*$ bounds the
second by $\sqrt{T_N^{(0)}T_N^{(2)}}$. These are precisely
\eqref{eq:tailerrors}; every omitted count is paid, and the retained
signed terms are unchanged.
\end{proof}

Since $h_*d_0=H_\lambda$, the finite analytical reconstruction from the
conserved initial count sum is
\begin{align}
 \widehat H_{\lambda,N}(k)&=H_{1/2}(k)-h_*\sum_{j=1}^Nd_j(t,k),
 \label{eq:tailreconstruction}\\
 \widehat H_{\lambda,N}'(k)&=H_{1/2}'(k)-h_*\sum_{j=1}^N\partial_kd_j(t,k).
 \notag
\end{align}
Its exact errors are $-h_*R_N$ and $-h_*\partial_kR_N$, respectively.
Thus the bounds in \eqref{eq:tailerrors}, multiplied by $h_*$, give
absolute errors for $H_\lambda$ and its first real-$k$ derivative. They
are uniform over all real $k$, not relative errors at zeros or collisions.

\subsection{Endpoints, score balance and the direct baseline}
At $t=0$, positive-count masses vanish, the quotient in
\eqref{eq:hcratio} is undefined, and the reconstruction is exact with
$Q_N=0$ and zero seeds. The normalized source expression has a limit
there, but is not a positive-count physical event. The endpoint $t=T$,
$\lambda=-1/2$, is included in both theorems. At $k=0$ the normalized
trace is exactly one; sign subsequences are only for fixed $k\ne0$ and
have no uniform finite starting index near zero. The phase estimate is
uniform on each compact real $k$ set; the tail estimates are uniform
absolute bounds for all real $k$. Neither gives an all-height relative
margin, bounded complex measurement, infinite-time result or uniform
$\gamma\downarrow0$ conclusion.

The one assessed source-score identity provides a useful comparison
\cite{TNDSourceCountDecision,TNDHighCountReview}. For $s(v)=\Phi'(v)/\Phi(v)$,
complete smooth reflection and the full first-derivative tails justify
integration of $(w_t(v)\sin(kv))'$ on the full line, with both boundary
terms zero. Since $w_t'=(2\lambda(t)v+s(v))w_t$ and
$C_k=-\E_t[X\sin(kX)]$, this gives
\begin{equation}\label{eq:scorebalance}
 kC(t,k)=2\lambda(t)C_k(t,k)-\E_t[s(X)\sin(kX)].
\end{equation}
At $C=C_k=0$ the correlation also vanishes: it is not a third independent
signed constraint. Subtracting a linear term from $s$ only changes the
$C_k$ coefficient; a further source-sensitive comparison remains unpaid.

The positive concentration proof applies to any source with the two-sided
comparison \eqref{eq:hccompare}; it does not characterize exact theta
arithmetic. The reflected control of Section~\ref{sec:control} fails its
canonical double-exponential upper bound, while retaining the generic
complete count and first-derivative reconstruction identities with its own
normalization. Its collision is compatible with all those identities.
No claim is made that the control lacks its own high-order oscillation.

Direct differentiation of the same scalar source gives
\eqref{eq:hcderivatives}; direct expansion of $\ee^{\alpha r^2}$ against
$\ee^{\lambda r^2}\Phi(r)$ gives exactly the same omitted kernel and
bounds as \eqref{eq:tailerrors}. There is no estimating or computational
gain over that baseline. The phase theorem excludes only eventual one-sign
literal count sectors at fixed nonzero real $k$; it leaves correlated
cancellations, phase-dependent patterns and finite-sector mechanisms open.
The tail theorem pays absolute remainders without giving a new tail sign
or collision restriction. Complete source/radial tails, every angular shell
and sign, the supplied loader and absolute calibration are unchanged.
Analytical finite-count reconstruction does not assume experimental access
to the specified cosine after nonzero counts.
